\section{支撑材料与结果生成链}

论文入口为\texttt{main.tex}，正文四至七章分别为四个问题，每章依次包含分析、模型建立和求解。原\texttt{example.tex}、类文件及样式文件保留。计算代码、原始模板和第一版源代码分别保存在\path{code/}、\path{originals/}、\path{versions/v1_before_settlement_revision/}；分问重排前的正文与代码另存\path{versions/v2_before_question_restructure/}。在项目目录使用\texttt{latexmk -xelatex main.tex}可编译当前论文。

\begin{table}[htbp]\centering\small
\caption{修订稿计算入口与结果文件}\label{tab:files}
\begin{tabularx}{\linewidth}{p{0.43\linewidth}X}\toprule
文件 & 用途\\\midrule
\path{code/data_forecast.py} & 原始附件读取及因果预测基础\\
\path{code/optimize_dispatch.py} & 保留原版LP、残差场景与规则执行\\
\path{code/dispatch_v2.py} & 修订控制器、跨日续接与逐段结算\\
\path{code/calibrate_v2.py} & 一月费用选型、匹配参数与冻结配置\\
\path{code/run_main_v2.py} & 四组未筛选母模型结果\\
\path{code/outer_candidate_selector.py} & 候选生成与固定执行规则回放评分\\
\path{code/run_outer_candidate_experiment.py} & 母模型与全时点筛选对照\\
\path{code/run_outer_timing_ablation.py} & 四次筛选时点消融与配置冻结\\
\path{code/publish_final_strategy.py} & 最终配置汇总、主结果图表及附件数据\\
\path{code/experiments_v2.py} & 场景、风险、备用、反馈及控制对照\\
\path{code/diagnose_execution_gap.py} & 原版全决策点规划与规则回放诊断\\
\path{code/diagnose_inner_v2.py} & 指定四日的同合同两控制器诊断\\
\path{code/time_index.py} & 物理时间、模板列序及分组映射\\
\path{code/make_paper_assets_v2.py} & 母模型数表、插图及原对照\\
\path{code/make_question_tables.py} & 分问编排中的预测与场景对照表\\
\path{code/make_visual_atlas.py} & 热力矩阵、瀑布图与场景阴影带\\
\path{code/scenario_convergence.py} & 真实场景缩减迭代记录与收敛分位带\\
\path{code/check_dispatch_v2.py} & 数值守恒、未来扰动及价格机制检查\\
\path{code/export_final_workbooks.mjs} & 从最终导出数据重填两本滚动工作簿\\
\path{code/verify_final_workbooks.py} & 五本最终附件的逐值、标签及汇总检查\\
\path{code/check_final_strategy_delivery.py} & 最终执行、图表、正文与附件独立核对\\
\path{code/run_revision.ps1} & 默认复用冻结配置的修订复算入口\\
\bottomrule\end{tabularx}
\end{table}

\texttt{selection\_v2.json}保存每个一月候选、费用、末库存、母模型配置及固定参数匹配值。\texttt{prepared\_v2\_*.npz}为已选预测输入，\texttt{*\_v2.npz}为334天母模型执行结果；旧\texttt{*\_full.npz}是原版365天记录，提供冻结的一月运行前缀。最终主结果中，问题二及第四问日前版仍读取母模型数组，问题三及第四问滚动版分别读取\path{outputs/outer_candidate_test/}中的\texttt{q3\_selected.npz}及\texttt{q4\_3\_selected.npz}。

筛选时点的四次回测、逐决策评分、配置及原始检查保存在\path{outputs/outer_timing_ablation/}；其中\texttt{final\_strategy.json}明确每问所选配置和执行数组路径。原运行清单保留当时输入的指纹，论文与附件随后获准更新，不回写旧清单伪装为原实验输入。当前发布文件及其来源另由\path{outputs/final_strategy/}中的清单记录。

每个主结果中，\texttt{plan}为零时合同，\texttt{adjusted}为每段最终有效常规量，\texttt{charge}、\texttt{discharge}为电池外侧实际电量，\texttt{energy}每一天含145个边界状态。母模型另以\texttt{reference\_energy}保存外层参考路径，最终两组滚动数组不包含该字段，故不报告混合含义的参考SOC偏差。\texttt{emergency}与\texttt{spill}分别为紧急与弃余电量；\texttt{plan\_cost}、\texttt{adjustment\_cost}及\texttt{emergency\_cost}之和为\texttt{total\_cost}。所有量保持完整精度，表内显示再按相应小数位舍入。

最终主执行数组汇入\path{outputs/final_strategy/results_export.json}，同目录的\texttt{summary.json}及\texttt{daily\_summary.csv}保存主指标、逐日费用和库存。工作簿与正文数表使用对应最终数组和统一索引，不从PDF抄录汇总数值。\path{data/results_export.json}保留母模型导出，不能用于覆盖最终滚动附件。原\path{experiments_v2.json/npz}及规划执行诊断仍保存母模型对照，不与新主案费用拼成未经运行的实验。

\section{时间映射与附件5填写}

附件5购电页保持原模板首标签“0:10--0:20”和末标签次日“0:00--0:10”。首末标签的原始文本在不同工作簿中略有差别，均逐字保留。以下列出源样本、物理模型与模板显示的对应，不据标签再平移数组：

\begin{center}\small
\begin{tabular}{llll}\toprule
数组索引 & 源样本右端点 & 物理时段 & 保留模板显示\\\midrule
$t=0$ & 00:10 & 00:00--00:10 & 0:10--0:20\\
$t=60$ & 10:10 & 10:00--10:10 & 10:10--10:20\\
$t=143$ & 次日00:00 & 23:50--24:00 & 次日0:00--0:10\\
\bottomrule\end{tabular}
\end{center}

具体原标签、Excel行列与四小时块均冻结于\path{data/time_mapping.json/csv}。问题一为纵向144行，其他购电页为横向144列；物理10:00--10:10从同一$t=60$取值。四小时储能区间依次汇总24段，紧急页将同一天相邻正紧急量合并，跨日不合并，无紧急量日期写“无”和0。源样本按右端点平均功率解释是明示建模假设，不将模板显示差异说成题目已给定的物理事实。

计划页的全天购电量及费用分别为$\sum_tg_t$及$\sum_tp_tg_t$；调整页填$a_t$而非$a_t-g_t$，全天费用填式\eqref{eq:settlement}的全部实际账单。没有调整页的问题二与问题四静态案，其完整账单由逐日CSV和正文给出。五本已填文件保存在\texttt{附件/附件5/}，空白原模板仅保存在\texttt{originals/附件5/}，交付时必须区别二者。

原母模型填报曾逐格核对288720个计划及调整数值，记录保留在\path{outputs/workbook_previews/}。本次按最终配置重写\texttt{result3.xlsx}及\texttt{result4-3.xlsx}，另外三本保持原方案；两本滚动附件的计划及调整区域共192384个数值，并重新核对储能聚合、紧急区间、日期与模板标签。当前文件指纹、导出与独立核对记录保存在\path{outputs/final_strategy/}，覆盖前文件另行备份。文件一致性检查不替代独立复现模型和人工论文审查。

\section{数值可行性与因果信息检查}

四个最终主案检查供给量不小于净需求；允许弃余时，使用扣除弃置量后的等式残差核对账本，不将优化约束改为强制全部电量使用。另检查SOC递推、日际连续性、1200--10800 kWh库存界、每段$5000/6$ kWh充放电上限及逐段互斥。计划费用始终按全部零时合同量结算，富余或未使用部分不抵扣。

母模型的未来扰动检查见\path{data/numerical_checks_v2.json}。新增候选模块还检查了逐场景回放与标量执行器一致性、日内保留当前合同、原计划结算基准、配对场景价格、终端偏离期望以及未观测实值和未发布预报扰动后已执行前缀不变；时点消融检查关闭筛选后仍求原场景LP。详细结果在两个外层实验目录的检查记录中。有限因果检查不等同于任意输入的形式证明，也不能使开发期选型变成独立测试。

以下节选显示外层交易与内层动作的职责。完整数值实现及参数保存在源文件中。

\begin{lstlisting}[language=Python,basicstyle=\ttfamily\footnotesize]
for t in range(144):
    if t in allowed_updates:
        if screening_enabled(question, t):
            candidates = candidate_plans(legal_forecast, observed_energy,
                                         original_day_ahead_plan, contract)
            contract = select_by_fixed_rule_replay(candidates,
                       same_historical_scenarios, observed_energy)
        else:
            contract = outer_plan(legal_forecast, observed_energy,
                                  original_day_ahead_plan)
    # Contract remains fixed inside the storage controller.
    action = controller(contract[t:], current_net, current_price,
                        latest_legal_future_forecast, observed_energy)
    execute_first_action(action)

plan_cost = sum(price * original_plan)
adjustment_cost = sum(price * (1.5 * maximum(final_plan-original_plan, 0)
                              + 0.5 * maximum(original_plan-final_plan, 0)))
emergency_cost = sum(5 * price * emergency)
total_cost = plan_cost + adjustment_cost + emergency_cost
\end{lstlisting}

\section{指定日期费用分解}

以下附表将正文指定日期的常规量及三项费用并列，便于与工作簿计划页、调整页及紧急页交叉定位。
\begin{table}[htbp]\centering\small
\caption{问题2指定日期计划、调整与费用分解}
\begin{tabularx}{\linewidth}{lRRRRR}\toprule
日期 & 计划量/kWh & 生效量/kWh & 计划费/元 & 调整费/元 & 紧急费/元\\\midrule
2025-03-20 & 67796.64 & 67796.64 & 41219.41 & 0.00 & 0.00\\
2025-06-21 & 40036.30 & 40036.30 & 23986.38 & 0.00 & 0.00\\
2025-09-23 & 69576.50 & 69576.50 & 43098.28 & 0.00 & 3285.98\\
2025-12-21 & 101008.53 & 101008.53 & 65454.44 & 0.00 & 0.00\\
\bottomrule\end{tabularx}\end{table}
